\subsection{Context and prior work}

The finite-vortex renormalized energy of Bethuel, Brezis, and
H\'elein~\cite{BethuelBrezisHelein1994} is an antecedent of the
infinite-system energy considered here. Sandier and Serfaty introduced
the planar functional in their analysis of microscopic vortex patterns
in the Ginzburg--Landau model. They proved existence of a minimum,
approximation of its value by square-periodic fields, and triangular
optimality among Bravais lattices~\cite[Theorems~1 and~2]{SandierSerfaty2012}.
Their Conjecture~1 asks whether the triangular value is also the minimum
over the full admissible class. Theorem~\ref{thm:main} resolves this
conjecture positively. The current normalization in that paper has one
point per area $2\pi/m$; our density-one convention corresponds to its
background $m=2\pi$.

The lattice comparison belongs to a longer history. The minimization of
convergent Epstein zeta sums at fixed covolume was developed by
Rankin~\cite{Rankin1953} and Cassels~\cite{Cassels1959}, with an error
noted in~\cite{Cassels1963Corrigendum} and a repair supplied by
Ennola~\cite{Ennola1964Lemma} and simplified by Diananda~\cite{Diananda1964}.
Montgomery proved triangular minimality of the lattice theta function
for every positive Gaussian parameter~\cite[Theorem~1]{Montgomery1988}.
These are comparisons within the lattice class; the issue here is the
lower bound for arbitrary admissible configurations.

For the confining-potential class in their Theorem~2, Sandier and
Serfaty show that microscopic limits of finite Coulomb minimizers
concentrate on minimizers of the planar renormalized
energy~\cite[Theorem~2]{SandierSerfaty2015}.  Its minimum also determines
a next-order coefficient for optimal logarithmic energy on the
sphere~\cite[Theorem~1.5]{BeterminSandier2018}.  Petrache and Serfaty
prove a conditional implication from universal optimality to minimality
for Coulomb and Riesz renormalized energies in their configuration-energy
normalization~\cite[Theorem~2 and Corollary~4]{PetracheSerfaty2020}.
Their periodic comparison uses tori compatible with the reference
lattice, not arbitrary square tori, and their planar conclusion is
conditional on the Cohn--Kumar conjecture.  The companion universal-optimality
theorem supplies the Gaussian comparison used in that
implication~\cite[Theorem~1.1 and Corollary~1.3]{UniversalOptimality2026}.
The present proof treats
the logarithmic functional directly and uses no universal-optimality
result as an input.

Two geometric precedents help explain the finite-torus argument.
Screening a point charge by a uniform disk yields a separation argument
for planar confined logarithmic minimizers; see
Lieb, Rougerie, and Yngvason~\cite[Lemma~3 and Appendix~A]{LiebRougerieYngvason2018}.
Voronoi moment inequalities combine sharp cell comparisons with Euler's
bound on the average number of sides, as in Gruber's proof of a theorem
of Fejes T\'oth~\cite{Gruber1999}. Affine bounds on cell energies, sharp
at the regular hexagon, also appear in the Wasserstein model of
Bourne, Peletier, and Theil for cells with at least three sides and mass
$v\ge1.5\cdot10^{-4}$~\cite[Lemma~12]{BournePeletierTheil2014}.
Here the cell comparison must arise from a trial potential that is
continuous off the sites and has the prescribed logarithmic
singularities. Its nonradial correction and
the two scalar estimates of Section~\ref{sec:calibration} address that
requirement, including sectors whose perpendicular foot lies outside
the corresponding edge.

The whole-plane reduction below uses Sandier and Serfaty's local energy
estimates and periodic approximation. The finite-torus argument
provides the comparison with the triangular value. Throughout, the
conclusion concerns energy per unit area: compact changes to a
configuration can preserve this value, so the minimum statement alone
does not classify all minimizing configurations.
